% ===================================================================
%  TABELLA nr. 12 — La staziun. --- Las viafierras. --- Las navs.
%  Traduction 2026 d'apres texto/io, controlee sur texto/fr, texto/ca
%  et texto/oc. Consigne : voir texto/rm/10-tabella-01.tex.
%
%  LE TABLEAU 11 A DONNE UN CRITERE ; LE TABLEAU 12 EST L'OCCASION DE
%  LE METTRE A L'EPREUVE, ET IL PASSE.
%
%  Le critere disait : une institution fondee sur place porte un nom
%  compose sur place ; une institution qui arrive du dehors deja
%  constituee porte le nom qui vient avec elle. On l'avait tire de
%  deux objets — la caisse d'epargne et les pompiers — et on l'avait
%  verifie deux fois, en luxembourgeois et en romanche.
%
%  Ce tableau-ci fournit un troisieme cas, et il n'etait pas prevu :
%  LE CHEMIN DE FER. Toutes les institutions des tableaux 1, 5, 6, 11
%  sont arrivees dans les vallees avec un nom etranger — l'ecole
%  normale, la banque, la prefecture, le theatre, la poste. Le chemin
%  de fer non. La ligne des Grisons a ete montee par une compagnie du
%  pays, avec de l'argent du pays, et elle porte depuis un nom
%  romanche : « Viafier Retica ». Le mot « viafier » lui-meme est un
%  compose fait a la maison — la voie et le fer —, non un emprunt.
%
%  CE N'EST DONC PAS QUE LES PETITES LANGUES SOIENT INCAPABLES DE
%  NOMMER LE MODERNE. Le chemin de fer est plus recent que la
%  prefecture, plus technique que la banque, plus etranger a la vie
%  d'une vallee qu'un theatre — et c'est lui qui a recu le nom
%  indigene, parce que ce sont les gens du lieu qui l'ont fait. LE
%  CRITERE NE REGARDE NI LA DATE NI LA TECHNIQUE : IL REGARDE QUI A
%  FONDE.
%
%  Et il devient predictif, ce qui est nouveau dans la serie : devant
%  une institution qu'on n'a pas encore examinee, on peut demander
%  qui l'a fondee avant de regarder comment elle se nomme. C'est le
%  premier tableau de la colonne ou la regle sert a prevoir et non a
%  decrire.
%
%  Le reste de la page est du materiel roulant, et il se partage
%  comme au tableau 10 : ce qui vient du sud garde le mot du sud
%  (« vagun », « locomotiva », « tender »), et ce qui a ete fait ici
%  se compose ici (« viafier », « chapstaziun », « stgaudapeis »).
% ===================================================================

%%K t12-apar-1 apar
\VUsaut{4.0mm}
\VUtitre{15.0pt}{60}{TABELLA nr. 12}
\VUsaut{2.4mm}
\VUpk{11.4pt}{40}{La staziun. --- Las viafierras. --- Las navs.}
\VUsaut{3.4mm}
\textsc{Nota.} --- \textit{Ils uraris duvrads en mintga naziun èn differents, nus
duvrain pia sco exempel las uras da la} \VUgras{tabella franzosa.}

%%K t12-01-1 p
1. --- Jau hai gist viagià tras la Germania. Avant la partenza hai jau cumprà in
\VUgras{cudeschet d'urari} \textsuperscript{(15)} per savair l'ura da partenza dals trens.
Suenter avair constatà ch'il tren il pli cumadaivel parta da Paris a las nov da
saira ed arriva a Strassburg a la una e tredesch minutas, hai jau preparà mes viadi,
ed il di suenter hai jau ma laschà manar a la staziun.

%%K t12-02-1 p
2. --- Cur ch'jau sun entrà en la gronda halla da partenza, davos in vegl
\VUgras{spiritual} \textsuperscript{(2)} da champagna, che purtava en maun ses \VUgras{satg da
viadi} \textsuperscript{(3)}, eran gia arrivads blers \VUgras{viagiaders}
\textsuperscript{(1)}.

%%K t12-02-2 p
Perquai ch'jau vuleva trametter in \VUgras{telegram} \textsuperscript{(5)}, hai jau laschà
mussar a mai d'in \VUgras{controllur} \textsuperscript{(6)} \VUgras{il biro dal telegraf}
\textsuperscript{(4)}.

%%K t12-03-1 p
3. --- Avant che passar en la \VUgras{sala da spetga} \textsuperscript{(7)} sun jau ì a ma
pervair d'in \VUgras{bigliet} \textsuperscript{(9)} cun return.

%%K t12-04-1 p
4. --- Jau n'hai betg spetgà ditg al \VUgras{guschet} \textsuperscript{(8)}, perquai ch'i era
mo pauca glieud là. In \VUgras{schuldà} \textsuperscript{(10)}, ch'era sin via en congedi, ha
gì l'amabladad da ma ceder sia plazza, uschia ch'jau hai gì avunda temp per dumandar
intginas infurmaziuns al \VUgras{chapstaziun} \textsuperscript{(13)}, ch'era en discurs cun
il \VUgras{commissari} \textsuperscript{(12)} da la surveglianza e cun il \VUgras{schandarm}
\textsuperscript{(11)} da servetsch.

%%K t12-tit-1 sub
\VUpk{10.2pt}{40}{Inscripziun da la bagascha.}
\VUsaut{3.0mm}

%%K t12-05-1 p
5. --- Suenter avair cumprà ina gasetta a la \VUgras{biblioteca} \textsuperscript{(14)}, hai
jau laschà pasar mia \VUgras{bagascha} \textsuperscript{(17)}, ch'in \VUgras{purtader}
\textsuperscript{(16)} aveva gist manà cun in \VUgras{char da pachets} \textsuperscript{(19)};
\VUgras{l'impiegà} \textsuperscript{(22)} incumbensà da la \VUgras{balanza} \textsuperscript{(21)}
m'ha fatg a savair che mia crappa peseva trentatschintg kilos; quai faseva in
\VUgras{surpais} da tschintg kilos, per il qual jau hai stuì pajar in supplement al
\VUgras{guschet} da la bagascha \textsuperscript{(23)}.

%%K t12-06-1 p
6. --- Perquai ch'jau era memia baud, hai jau gì temp d'observar la circulaziun dals
viagiaders en la staziun. Tscherts dumbravan giu u reprendevan lur bagascha pitschna
al \VUgras{deposit da bagascha} \textsuperscript{(25)}; auters consultavan las \VUgras{tabellas
d'urari} \textsuperscript{(26)}. Ils viagiaders ch'arrivavan sa preschevan a la
\VUgras{sortida} \textsuperscript{(32)} cun lur \VUgras{valisch} \textsuperscript{(24)} en maun.
Intgins spetgavan a la \VUgras{spartida da la bagascha} \textsuperscript{(27)} e
preschentavan lur \VUgras{bulletin} \textsuperscript{(29)} per reprender lur crappas, lur
\VUgras{chaschas} \textsuperscript{(28)} u lur \VUgras{chanasters} \textsuperscript{(30)}.

%%K t12-07-1 p
7. --- Ils \VUgras{impiegads da l'accisa} \textsuperscript{(33)} exploravan premurusamain
ils pachets per constatar sch'ins aveva insatge da declarar.

%%K t12-08-1 p
8. --- Tscherts viagiaders èn ids al \VUgras{buffet} \textsuperscript{(34)}, nua ch'in
\VUgras{camerier} \textsuperscript{(35)} sa dava fadia da als servir. Els ston sa prescher,
perquai ch'il \VUgras{tren} \textsuperscript{(36)}, ch'è in express, na resta betg ditg en
la staziun.

%%K t12-09-1 p
9. --- Lura sun jau ì al perron da partenza ed hai tschertgà in \VUgras{vagun}
\textsuperscript{(37)} direct, per betg esser sforzà da midar vagun durant il viadi. Jau
sun muntà en in vagun cun corridor, che cuntegna in \VUgras{cumpartiment}
\textsuperscript{(38)} per fimaders ed in cumpartiment resalvà per «dunnas sulettas».

%%K t12-tit-2 sub
\VUpk{10.2pt}{40}{Partenza da notg.}
\VUsaut{3.0mm}

%%K t12-10-1 p
10. --- Suenter avair muntà il \VUgras{stgalim} \textsuperscript{(40)}, serrà la
\VUgras{porta} \textsuperscript{(39)}, rullà si il \VUgras{storen}
\textsuperscript{(41)} e mess mia bagascha pitschna sin la \VUgras{rait}
\textsuperscript{(43)}, sun jau sesì sin il
\VUgras{banc} \textsuperscript{(42)}. Vesend il \VUgras{lampist} \textsuperscript{(44)} emprender
las lampas, hai jau patratgà ch'jau stoppia passentar la notg en il vagun ed hai
clamà l'impiegà incumbensà da la locaziun dals \VUgras{chaschins da chau}
\textsuperscript{(45)} per al dumandar in.

%%K t12-11-1 p
11. --- Il conductur è vegnì a verifitgar ils biglietts. Jau al hai dumandà pertge
ch'nus n'essan anc betg partids, perquai ch'i è l'ura endretga. El m'ha respondì ch'il
\VUgras{chaptren} \textsuperscript{(46)} è occupà cun far metter velos en il \VUgras{furgun}
\textsuperscript{(47)}. Plinavant n'aveva ins anc betg midà tut ils stgaudapeis
\textsuperscript{(48)}.

%%K t12-12-1 p
12. --- Ma nus vinsan prest a partir, perquai ch'il \VUgras{chauffeur}
\textsuperscript{(50)}, sin ses \VUgras{tender} \textsuperscript{(49)}, prendeva carvun per
excitar il fieu da la \VUgras{locomotiva} \textsuperscript{(54)}, e \VUgras{il maschinist}
\textsuperscript{(51)} spetgava mo il signal dal \VUgras{sutchapstaziun}
\textsuperscript{(52)}, ch'era sin il \VUgras{perron} \textsuperscript{(53)}.

%%K t12-13-1 p
13. --- La \VUgras{caldera} \textsuperscript{(56)} da la locomotiva grivlava suravia; il
\VUgras{chamin} \textsuperscript{(55)} vomitava torrents da fim mischedà cun \VUgras{vapur}
\textsuperscript{(60)}. In \VUgras{sivel} \textsuperscript{(59)} stridulant ha sunà ed ils
\VUgras{pistuns} \textsuperscript{(58)} han cumenzà a sa mover, impulsond las \VUgras{rodas}
\textsuperscript{(57)} da la pesanta maschina.

%%K t12-tit-3 sub
\VUsaut{3.0mm}
\VUpk{10.2pt}{40}{Partenza!}
\VUsaut{3.0mm}

%%K t12-14-1 p
14. --- La sveltezza dal tren, che curriva sin ils \VUgras{binaris} \textsuperscript{(64)},
è vegnida accelerada; ils \VUgras{perruns} \textsuperscript{(65)} èn spariids; nus avain
passà ils \VUgras{necessaris} \textsuperscript{(66)}, ed era vaguns garads sin l'auter
\VUgras{binari} \textsuperscript{(63)}: \VUgras{vaguns da durmir} \textsuperscript{(67)},
\VUgras{vagun-restaurant} \textsuperscript{(68)}, \VUgras{vagun da posta} \textsuperscript{(69)}.

%%K t12-noto-1 noto
\VUnotes{143.2mm}{%
\textit{(*) Nota.} --- \textit{S'enchapescha ch'la descripziun dal viadi è tenor la
frontiera d'avant la guerra.}%
}

%%K t12-15-1 p
15. --- Lura è la \VUgras{staziun da la rauba} \textsuperscript{(71)} passada avant noss
egls. Sut ils hangars chargiavan impiegads la \VUgras{rauba} \textsuperscript{(72)}
spedida cun trens plaun. \VUgras{Squadras} \textsuperscript{(76)} sper il \VUgras{reservoir
d'aua} \textsuperscript{(73)} manovravan \VUgras{vaguns da bes-chas} \textsuperscript{(75)} e
\VUgras{vaguns plats} \textsuperscript{(79)} per furmar in \VUgras{tren da rauba}
\textsuperscript{(80)}.

%%K t12-16-1 p
16. --- Nus avain traversà l'ultim \VUgras{sviamaint} \textsuperscript{(78)}, sper il qual
stueva il manader dals sviamaints; nus avain passà il \VUgras{disc da signal}
\textsuperscript{(86)} e la staziun è sparida dalunsch.

%%K t12-17-1 p
17. --- Prest avain nus traversà ina \VUgras{punt da fier} \textsuperscript{(74)} che passa
sur il \VUgras{flum} \textsuperscript{(81)}. Suenter avair traversà quella punt, avain nus
suandà il flum, sin il qual pussants \VUgras{remorchaders} \textsuperscript{(82)} tiravan
pesantas \VUgras{barcas da charg} \textsuperscript{(83)}. Nus avain plinavant vis ina
\VUgras{nav da passagiers} \textsuperscript{(84)}, cur ch'ella abordava in \VUgras{pontun}
\textsuperscript{(85)}.

%%K t12-18-1 p
18. --- Suenter avair passà a gronda sveltezza pliras staziuns, avain nus traversà
intgins \VUgras{tunnels} \textsuperscript{(87)} ed avain laschà davos nus
\VUgras{chasettas} nundumbraivlas \textsuperscript{(88)} da \VUgras{guardgias da
barriera}. Cunà dal tschitschim dal tren,
sun jau adurmentà profundamain.

%%K t12-tit-4 sub
\VUsaut{3.0mm}
\VUpk{10.2pt}{40}{Staziun da Deutsch-Avricourt.}
\VUsaut{3.0mm}

%%K t12-19-1 p
19. --- Jau sun subit vegnì dasdà da la vusch d'in impiegà che clamava:
«Deutsch-Avricourt! \textsuperscript{(*)}. Tuts giu!» I ha gì lieu la controlla da la
duana e tut las formalitads fastidiusas da la frontiera. Jau hai stuì adattar mia ura
da giaglioffa a l'\VUgras{ura gronda} \textsuperscript{(89)} da la staziun, ch'era
tschuncantatschintg minutas memia baud; apaina dasdà hai jau distinguì cun fadia il
\VUgras{grond punctader} \textsuperscript{(90)} dal \VUgras{pitschen} \textsuperscript{(91)}; il
\VUgras{quadrant} \textsuperscript{(92)} sez era dal tut confus pervi da la nebla da la
damaun.

%%K t12-20-1 p
20. --- Entant ch'ils duaniers fasevan lur controlla, hai jau, sbadagliond, deschifrà
ils \VUgras{affischs} \textsuperscript{(93)} ch'ornan ils mirs da la staziun. Jau hai
gentà per passentar il temp, hai consultà la charta da la \VUgras{rait}
\textsuperscript{(94)} tudestga per vesair quant lunsch ch'jau era anc da Strassburg, e
sun turnà en mes vagun, rinforzà da la speranza da cuntanscher prest il scopo da mes
viadi.

\VUsaut{6.0mm}
\VUfilet{22mm}
